\subsection{A finite section bound from one fibre}
\label{sec:finite-section-jets}

The upper section-rank estimate has a short geometric proof. It does
not require a uniform threshold in an asymptotic Hilbert--Samuel theorem.

\begin{proposition}\label{prop:finite-section-jets}
Let $X$ be an integral projective variety of dimension $D\ge1$ over
an algebraically closed field of characteristic zero. Let $A$ be nef,
and suppose a finite linear system in $H^0(X,A)$ gives a birational
map onto a projective image of degree $\delta>0$. Let $B$ be a line
bundle and $n\ge0$ an integer such that $nA-B$ is nef. For an integer $a\ge0$,
\[
 nA^D<(a+1)\delta
 \quad\Longrightarrow\quad
 h^0(X,B)\le\delta\binom{a+D}{D}.
\]
\end{proposition}

\begin{proof}
Blow up the base ideal, obtaining $\pi:Y\to X$ and an effective
Cartier divisor $E$ such that $L=\pi^*A-E$ is generated by the chosen
sections. Then $Y$ is integral and projective, $L$ is nef, and
$L^D=\delta$. A general linear projection of the image gives
$f:Y\to\mathbf P^D$ with $L=f^*\mathcal O(1)$, generically finite
of degree $\delta$. Choose $q$ in its finite \emph{\'etale} locus.
The fibre $Q=f^{-1}(q)$ consists of $\delta$ regular points. We claim
that the pullback of every section of $B$ is determined by its jets
through order $a$ at $Q$.

Suppose a nonzero section $s$ has all these jets zero. Its pullback
has an effective zero divisor $Z_s$. The nonfinite-fibre image of
$f$ has codimension at least two: positive-dimensional fibres over
a divisor would force a proper closed subset of the integral
$D$-fold $Y$ to have dimension at least $D$. A general line $\ell$
through $q$ therefore avoids that image and is not contained in
$f(Z_s)$. For $D=1$ take $\ell=\mathbf P^1$. The curve cycle
$C=f^{-1}(\ell)$ has class $L^{D-1}$. Every curve component dominates
$\ell$, hence none lies in $Z_s$; near $Q$, the curve is smooth
with multiplicity one. Each point of $Q$ contributes at least
$a+1$ zeros to $s|_C$. Consequently,
\[
 (a+1)\delta\le(\pi^*B)L^{D-1}
       \le n(\pi^*A)L^{D-1}\le n(\pi^*A)^D=nA^D,
\]
a contradiction. The middle inequalities use nefness. For the
last one, expand $(\pi^*A)^D-(\pi^*A)L^{D-1}$ as
$E(\pi^*A)\sum_{j=0}^{D-2}(\pi^*A)^{D-2-j}L^j$ if $D\ge2$;
each summand is nonnegative. For $D=1$ it is an equality.

Thus evaluation embeds $H^0(X,B)$ in the direct sum of the local
quotients $(\pi^*B)_z\otimes\mathcal O_{Y,z}/\mathfrak m_z^{a+1}$,
for $z\in Q$. Their dimensions are $\binom{a+D}{D}$, by counting
monomials in regular local parameters. This proves the bound.
The blowup, codimension-one finiteness and intersection facts used
here are recorded in the Stacks Project, tags
\href{https://stacks.math.columbia.edu/tag/01OF}{01OF},
\href{https://stacks.math.columbia.edu/tag/0AB5}{0AB5} and
\href{https://stacks.math.columbia.edu/tag/0BEL}{0BEL}.
\end{proof}

Suppose now that $t>D$ is an integer, $0\le\rho<1$ and
$\delta\ge(1-\rho)A^D$. Set
\[
 n=t+D+3,\qquad
 a=\left\lfloor\frac{t+D+3}{1-\rho}\right\rfloor,\qquad
 w=\delta\binom{t-1}{D},\qquad
 \eta=1-\frac{\binom{t-1}{D}}{\binom{a+D}{D}}.
\]
The proposition gives $(1-\eta)r\le w$ for any section subspace
of dimension $r$. The ratio defining $1-\eta$ is
$\prod_{j=1}^D(t-j)/(a+j)$. Its $j$th factor is at least $1-e_j$,
where
\[
 1-e_j=\frac{(1-\rho)(t-j)}{t+D+3+j(1-\rho)},\qquad
 0\le e_j<1,\qquad
 e_j\le\rho+\frac{D+3+2j}{t}.
\]
Using $1-\prod_j(1-e_j)\le\sum_j e_j$ gives
\[
 \boxed{\eta\le D\rho+\frac{2D^2+4D}{t}
               \le D\rho+\frac{4D^2}{t}\quad(D\ge2).}
\]
If a separate lower estimate gives $w\le r$ and the first majorant
$\Theta$ is less than one, the actual excess rank satisfies
\[
             0\le\frac{r-w}{w}\le\frac{\eta}{1-\eta}
                                  \le\frac{\Theta}{1-\Theta}.
\]
This is a relative rank comparison, not a bound for an arithmetic
degree. Its application retains the original projective model,
section realization, nefness, birationality and degree comparison.
The signed cusp mean and the other arithmetic section remainders
remain separate requirements.
